% 原书物理页=印刷页31--35上（止于第3章章名前）；2.5 全节.
\section{半单 Lie 代数}\label{sec:ssl-2-5}
若 $\mathfrak h,\mathfrak h'$ 是 $K$-Lie 代数 $\mathfrak g$ 的两个可解理想，则它们的和 $\mathfrak h+\mathfrak h'$ 仍可解，因为它是直和 $\mathfrak h\oplus\mathfrak h'$ 的一个商.\footnote{译注：原文的这一理由照录. 若指 Lie 代数的商，则一般不成立，因为两理想未必互相交换，求和映射未必是 Lie 代数同态. 可解性的结论可由 $(\mathfrak h+\mathfrak h^{\prime})/\mathfrak h\simeq\mathfrak h^{\prime}/(\mathfrak h\cap\mathfrak h^{\prime})$ 及可解代数对扩张封闭得到.} 据此可以定义\newterm{ssl-radical}{根基}{radical} $\operatorname{rad}(\mathfrak g)$，即 $\mathfrak g$ 的最大可解理想.
\begin{definition}\label{def:ssl-2-5-1}
$K$-Lie 代数 $\mathfrak g$ 称为
\begin{itemize}[label=\textendash]
\item \newterm{ssl-simple}{单}{simple}的，如果 $\dim(\mathfrak g)>1$ 且它仅有的理想是 $\{0\}$ 与 $\mathfrak g$；
\item \newterm{ssl-semisimple}{半单}{semisimple}的，如果 $\operatorname{rad}(\mathfrak g)=0$（亦即 $\mathfrak g$ 没有非零的 Abel 理想）；
\item \newterm{ssl-reductive}{约化}{reductive}的，如果 $\operatorname{rad}(\mathfrak g)=\mathfrak z(\mathfrak g)$.
\end{itemize}
\end{definition}
\begin{remark*}
``单''的定义中的条件 $\dim(\mathfrak g)>1$ 是为了排除``平凡''代数 $\mathfrak g=K$，它的仅有的理想同样是 $\{0\}$ 与 $\mathfrak g$. 在这些定义下，有
\[\mathfrak g\ \text{单}\Longrightarrow\mathfrak g\ \text{半单}\Longrightarrow\mathfrak g\ \text{约化}\]
\end{remark*}
\begin{exercise*}
\begin{enumerate}[label=\roman*)]
\item 证明 $\mathfrak g$ 单 $\Rightarrow$ $\mathfrak g$ 半单.
\item 证明对任意的 $\mathfrak g$，$\mathfrak g/\operatorname{rad}(\mathfrak g)$ 是半单的.
\item 证明半单代数的直和是半单的.
\end{enumerate}
\end{exercise*}
我们将看到，所有经典 Lie 代数都是约化的，而且只要其中心平凡，它们甚至是单的.\footnote{译注：此处原文的概括未排除小维数例外. 例如 $\mathfrak{so}_4(\mathbb C)\simeq\mathfrak{sl}_2(\mathbb C)\oplus\mathfrak{sl}_2(\mathbb C)$ 中心为零，却不是单 Lie 代数；正文照录.}

\begin{numpara}[借助 Killing 型的刻画]\label{np:ssl-2-5-2}
从现在起假设 $K$ 是特征为零的代数闭域.
\begin{lemma*}
设 $(V,r)$ 是 $\mathfrak g$ 的一个不可约表示. 则每个 $X\in\operatorname{rad}(\mathfrak g)$ 都以标量作用，从而 $r([\mathfrak g,\operatorname{rad}(\mathfrak g)])=0$.
\end{lemma*}
\begin{proof}
由 Lie 定理，$V$ 中存在 $\operatorname{rad}(\mathfrak g)$ 的特征向量. 记与之相伴的特征值为 $\lambda\colon\operatorname{rad}(\mathfrak g)\longrightarrow K$，与之相伴的特征子空间为 $V_{\lambda}$. 我们只须证明 $V_{\lambda}$ 在 $\mathfrak g$ 下稳定. 而对一切 $Y\in\mathfrak g$、$X\in\operatorname{rad}(\mathfrak g)$ 与 $v\in V_{\lambda}$，有 $XYv=[X,Y]v+YXv=\lambda([X,Y])v+\lambda(X)Yv$. 与 Lie 定理的证明中一样，可证 $\lambda([X,Y])=0$，由此即得结论.
\end{proof}
\begin{corollary*}
若 $\mathfrak g$ 有一个表示 $(V,r)$，使得双线性型 $B_r$ 非退化，则 $\mathfrak g$ 是约化的.
\end{corollary*}
\begin{proof}
由引理，理想 $[\mathfrak g,\operatorname{rad}(\mathfrak g)]$ 在 $\mathfrak g$ 的每个不可约表示上都平凡作用. 取 $0\subset V_1\subset V_2\subset\cdots\subset V_l=V$ 为 $V$ 的一条递增子表示链，其逐次商 $V_i/V_{i-1}$ 都不可约. 于是对一切 $Y\in\mathfrak g$ 有 $r(Y)(V_i)\subset V_i$，而对 $X\in[\mathfrak g,\operatorname{rad}(\mathfrak g)]$ 有 $r(X)(V_i)\subset V_{i-1}$. 因此 $r(X)r(Y)$ 幂零，从而迹为零；由此推出 $[\mathfrak g,\operatorname{rad}(\mathfrak g)]$ 含于 $B_r$ 的核. 由于 $B_r$ 已假设为非退化，得 $[\mathfrak g,\operatorname{rad}(\mathfrak g)]=0$，这等价于 $\operatorname{rad}(\mathfrak g)\subset\mathfrak z(\mathfrak g)$. 反向包含是同义反复.
\end{proof}
\begin{example*}
把这个推论应用于经典 Lie 代数的``标准''表示，可以看出 $\mathfrak{gl}_n(K)$ 与 $\mathfrak u(n)$（当 $K=\mathbb{R}$）是约化的，而 $\mathfrak{sl}_n(K)$、$\mathfrak{so}_n(K)$（当 $n>2$）、$\mathfrak{sp}_{2n}(K)$、$\mathfrak{su}(n)$（当 $K=\mathbb{R}$）是约化的，甚至是半单的，因为（习题）它们的中心为零.
\end{example*}
\begin{theorem*}[Cartan]
$K$-Lie 代数 $\mathfrak g$ 是半单的，当且仅当 Killing 型 $B_{\operatorname{ad}}$ 非退化.
\end{theorem*}
\begin{proof}
设 $B_{\operatorname{ad}}$ 非退化. 这显然蕴涵 $\operatorname{Ker}(\operatorname{ad})=\mathfrak z(\mathfrak g)=0$，而由上面的推论，它也蕴涵 $\mathfrak g$ 是约化的. 于是 $\mathfrak g$ 半单.

反之，设 $\mathfrak g$ 半单. 则 $\mathfrak z(\mathfrak g)=0$，故 $\operatorname{ad}\colon\mathfrak g\longrightarrow\mathfrak{gl}(\mathfrak g)$ 是单射. 令 $\mathfrak h=\mathfrak g^{\perp}$，即 $B_{\operatorname{ad}}$ 的核. 由 Cartan 判据，$\mathfrak h$ 可解，故 $\mathfrak h=0$.
\end{proof}
\begin{theorem*}
半单 $K$-Lie 代数 $\mathfrak g$ 是单理想的直和. 特别地，$\mathfrak g=[\mathfrak g,\mathfrak g]$.
\end{theorem*}
\begin{proof}
设 $\mathfrak h$ 是 $\mathfrak g$ 的一个理想，$\mathfrak h^{\perp}$ 是它关于 $\mathfrak g$ 的 Killing 型的正交补. 那么 $\mathfrak h\cap\mathfrak h^{\perp}$ 的 Killing 型（它是 $\mathfrak g$ 的 Killing 型的限制）为零，故 $\mathfrak h\cap\mathfrak h^{\perp}$ 是可解理想，从而为零. 由于 $B_{\operatorname{ad}}$ 非退化，由此推出 $\mathfrak g=\mathfrak h\oplus\mathfrak h^{\perp}$，并且 $\mathfrak h$ 与 $\mathfrak h^{\perp}$ 的 Killing 型都非退化. 对维数作归纳即可完成证明.
\end{proof}
\end{numpara}

\begin{numpara}[借助表示的刻画]\label{np:ssl-2-5-3}
仍设 $K$ 是特征为零的代数闭域.
\begin{theorem*}[Weyl]
$K$-Lie 代数 $\mathfrak g$ 是半单的，当且仅当它的所有（有限维）表示都完全可约.
\end{theorem*}
Weyl（在 $K=\mathbb{C}$ 时）对这条定理的原始证明使用了他的``酉技巧''（unitary trick）：找到 $\mathfrak g$ 的一个可积成紧 Lie 群的\newterm{ssl-real-form}{实形式}{real form}，以便利用众所周知的事实——紧群的连续有限维表示都完全可约（这得益于 Haar 测度的存在性）. 参见 [GALM2]，特别是其中的推论 1.12.7 及其``应用''. 下面给出一个更直接的证明，它出自 Casimir（一位物理学家！）.
\begin{proof}
若诸表示都完全可约，则特别地伴随表示完全可约. 把 $\mathfrak g$ 写成 $\mathfrak g=\bigoplus_i\mathfrak g_i$，其中诸 $\mathfrak g_i$ 是 $(\mathfrak g,\operatorname{ad})$ 的不可约子表示. 于是诸 $\mathfrak g_i$ 都是理想（因为它们在 $\operatorname{ad}_{\mathfrak g}$ 下稳定），并且是单的或一维的（因为不可约）. 而 $\mathfrak g$ 不可能有一维的直和因子，因为我们已经看到 $K$ 的表示不一定完全可约. 故 $\mathfrak g$ 半单.

现在设 $\mathfrak g$ 半单. 设 $(V,r)\in\operatorname{Rep}_K(\mathfrak g)$，$W\subset V$ 是一个子表示；我们要找 $W$ 的一个在 $\mathfrak g$ 下稳定的补. 为此，只须找到一个 $\varphi\in\operatorname{Hom}_{K\mathfrak g}(V,W)$\footnote{译注：原文此处下标作 $K\mathfrak g$，此处保留. 按证明中所用的意义，这是 $\mathfrak g$-模之间的同态，也即第~\ref{sec:ssl-2-2}~节所述的 $\mathcal U\mathfrak g$-模同态.}，使得 $\varphi|_W=\lambda\,\mathrm{id}_W$ 且 $\lambda\neq 0$. 因为此时将有 $V=W\oplus\operatorname{Ker}(\varphi)$，且 $\operatorname{Ker}(\varphi)$ 稳定.

\emph{第一步：}化归为 $V/W$ 是单位表示的情形. 考察内部 $\mathrm{Hom}$ 表示 $H=\mathcal{H}om_K(V,W)$ 及其子空间 $A:=\{\varphi\in H,\ \varphi|_W\in K\,\mathrm{id}_W\}$ 与 $B:=\{\varphi\in A,\ \varphi|_W=0\}$. 若 $\varphi\in A$，则对一切 $X\in\mathfrak g$ 有 $(X\varphi)(w)=X\varphi(w)-\varphi(Xw)=X\lambda w-\lambda Xw=0$，故 $X\varphi\in B$. 特别地，$A$ 与 $B$ 是 $H$ 的子表示，且 $B/A$ 是单位表示.\footnote{译注：原文商写作 $B/A$，此处照录；由 $B\subset A$ 及随后寻找 $B$ 在 $A$ 中的补空间可知，应为 $A/B$.} 假设 $B$ 有一个补，即一条直线 $D\subset A^{\mathfrak g}$ 使 $B\oplus D=A$. 那么，如上面所解释的，就有 $W\oplus\operatorname{Ker}(\varphi)=V$.

\emph{第二步：}Casimir 算子与归纳. 从现在起设 $V/W$ 是单位表示，并对 $\dim(V)$ 作归纳，证明 $W$ 的稳定补直线的存在性. 固定一条由子表示构成的递增链 $0=V_0\subset V_1\subset\cdots\subset V_{l-1}=W\subset V_l=V$，使得对一切 $i=1,\cdots,l$，$V_i/V_{i-1}$ 都不可约.
若所有 $V_i/V_{i-1}$ 都同构于单位表示，那么我们就得到了 $V$ 的一个完全旗 $\mathcal V$，使 $r(\mathfrak g)\subset\mathfrak b^{\mathrm{nilp}}(\mathcal V)$. 但由于对一切 $n$ 有 $\mathfrak g=[\mathfrak g,\mathfrak g]=C^n\mathfrak g$，取 $n=\dim(V)$ 便得 $r(\mathfrak g)=r(C^n\mathfrak g)=C^n\mathfrak b^{\mathrm{nilp}}(\mathcal V)=0$，于是 $V$ 是平凡表示，$W$ 的任何补直线都稳定.
否则，存在 $i$ 使 $(U,r_U):=V_i/V_{i-1}$ 不同于单位表示. 此时 $r_U(\mathfrak g)$ 非零，故可取 $\mathfrak g$ 的一个\newterm{ssl-simple-factor}{单因子}{simple factor} $\mathfrak g_U$ 使 $r_U(\mathfrak g_U)\neq 0$. 于是 $\mathfrak g_U\xrightarrow{\ \sim\ }r_U(\mathfrak g_U)$，因而 $r_U(\mathfrak g_U)$ 是单的，且与 $r_U$ 相伴的双线性型 $B_{r_U}$ 非零（否则 Cartan 判据将蕴涵 $r_U(\mathfrak g_U)$ 可解）. 它的核是 $\mathfrak g_U$ 的真理想，故为零；最终，$B_{r_U}$ 在 $\mathfrak g_U$ 上非退化. 令 $C_U$ 为相应的 Casimir 算子，它属于 $\mathcal{U}\mathfrak g_U$ 的中心，从而属于 $\mathcal{U}\mathfrak g$ 的中心（因为 $\mathcal{U}(\mathfrak h\oplus\mathfrak h')=\mathcal{U}\mathfrak h\otimes_K\mathcal{U}\mathfrak h'$）. 于是 $r(C_U)$ 是 $V$ 上与 $\mathfrak g$ 的作用交换的自同态，故 $\mathfrak g$ 使它的诸广义特征空间稳定. 我们有 $r(C_U)(V)\subset W$，因此 $0$ 是特征值，且广义特征空间 $V_0:=\operatorname{Ker}(r(C_U)^n)$ 不含于 $W$. 令 $V^0$ 为 $r(C_U)$ 的其余广义特征空间之和. 于是得到 $V$ 到稳定子空间的一个直和分解 $V^0\oplus V_0$. 若能证明 $V^0\neq 0$，就可以对 $V_0$ 应用归纳假设. 因此我们必须证明 $r(C_U)$ 有一个非零特征值. 为此只须验证 $r_U(C_U)$ 有一个非零特征值. 注意由 Schur 引理，$r_U(C_U)$ 是一个位似. 于是只须证明它的迹非零. 而 $\operatorname{tr}(r_U(C_U))=\sum_{i=1}^{\dim\mathfrak g_U}\operatorname{tr}(r_U(X_i)r_U(X_i^{*}))$，其中 $(X_i)$ 是 $\mathfrak g_U$ 的一组基，$(X_i^{*})$ 是它关于双线性型 $B_{r_U}(X,Y)=\operatorname{tr}(r_U(X)r_U(Y))$ 的对偶基. 于是 $\operatorname{tr}(r_U(C_U))=\dim\mathfrak g_U\neq 0$.
\end{proof}
\begin{corollary*}
$K$-Lie 代数 $\mathfrak g$ 是约化的，当且仅当 $\mathfrak g=\mathfrak z(\mathfrak g)\oplus[\mathfrak g,\mathfrak g]$ 且 $[\mathfrak g,\mathfrak g]$ 半单.
\end{corollary*}
\begin{proof}
条件的充分性是明显的. 证其必要性. 设 $\mathfrak g$ 约化. 于是伴随表示 $\operatorname{ad}\colon\mathfrak g\longrightarrow\mathfrak{gl}(\mathfrak g)$ 通过 $\operatorname{ad}(\mathfrak g)=\mathfrak g/\mathfrak z(\mathfrak g)$ 分解，而后者是半单的. 故它完全可约. 由于 $\mathfrak z(\mathfrak g)\subset\mathfrak g$ 在 $\operatorname{ad}(\mathfrak g)$ 下稳定，存在分解 $\mathfrak g=\mathfrak z(\mathfrak g)\oplus\mathfrak g_1\oplus\cdots\oplus\mathfrak g_k$，其中诸 $\mathfrak g_i$ 不可约，从而单. 由此 $[\mathfrak g,\mathfrak g]=[\mathfrak g_1,\mathfrak g_1]\oplus\cdots\oplus[\mathfrak g_k,\mathfrak g_k]=\mathfrak g_1\oplus\cdots\oplus\mathfrak g_k$.
\end{proof}
\end{numpara}

\begin{numpara}[Jordan 分解]\label{np:ssl-2-5-4}
下面是 Weyl 完全可约性定理的第一个应用.
\begin{theorem*}
设 $X$ 是半单 Lie 代数 $\mathfrak g$ 的一个元素. 存在唯一的分解 $X=X_s+X_n$，使得对任何表示 $(r,V)$，都有 $r(X_s)=r(X)_s$ 与 $r(X_n)=r(X)_n$，其中 $r(X)_s+r(X)_n$ 是 $V$ 上自同态 $r(X)$ 的 Jordan 分解. 此外，有 $[X_s,X_n]=0$.
\end{theorem*}
\begin{proof}
\emph{第一步：}对任何表示 $(V,r)$，有 $r(X)_s\in r(\mathfrak g)$ 且 $r(X)_n\in r(\mathfrak g)$.

首先注意，由于 $r(X)_s$ 与 $r(X)_n$ 是 $r(X)$ 的多项式，$V$ 的任何子表示 $W$ 都在 $r(X)_s$ 与 $r(X)_n$ 下稳定. 此外，由于 $r(X)_n$ 幂零，有 $\operatorname{tr}(r(X)_n|_W)=0$；又由于 $\mathfrak g=[\mathfrak g,\mathfrak g]$，还有 $\operatorname{tr}(r(X)_s|_W)=\operatorname{tr}(r(X)|_W)=0$.

另一方面，由于 $\operatorname{ad}_{r(X)_s}$ 与 $\operatorname{ad}_{r(X)_n}$ 是 $\operatorname{ad}_{r(X)}$ 的多项式（回忆 $\operatorname{ad}_{r(X)_s}+\operatorname{ad}_{r(X)_n}$ 是 $\operatorname{ad}_{r(X)}$ 的 Jordan 分解），并且 $[r(X),r(\mathfrak g)]=\operatorname{ad}_{r(X)}(\mathfrak g)\subset\mathfrak g$\footnote{译注：原文此式及下面两式中像集均写作 $\mathfrak g$；按上下文（$\mathfrak g$ 经 $r$ 映入 $\mathfrak{gl}(V)$，且下文以 $[u,r(\mathfrak g)]\subset r(\mathfrak g)$ 定义 $\mathfrak n$）应为 $r(\mathfrak g)$.}，同样有 $[r(X)_s,r(\mathfrak g)]=\operatorname{ad}_{r(X)_s}(\mathfrak g)\subset\mathfrak g$ 与 $[r(X)_n,r(\mathfrak g)]=\operatorname{ad}_{r(X)_n}(\mathfrak g)\subset\mathfrak g$. 于是引入\newterm{ssl-normalizer}{规范化子}{normalizer} $\mathfrak n:=\{u\in\mathfrak{gl}(V),\ [u,r(\mathfrak g)]\subset r(\mathfrak g)\}$，它是 $\mathfrak{gl}(V)$ 的一个 Lie 子代数，把 $r(\mathfrak g)$ 作为一个理想包含在内；再引入它的 Lie 子代数 $\mathfrak n':=\{u\in\mathfrak n,\ \text{对每个 }\mathfrak g\text{-子模 }W\subset V,\ u(W)=W\ \text{且}\ \operatorname{tr}(u|_W)=0\}$\footnote{译注：原文作 $u(W)=W$；按文意应为 $u(W)\subset W$（只需 $u$ 保持 $W$ 稳定，否则如下文所言 $r(\mathfrak g)$ 及 $r(X)_s$、$r(X)_n$ 不必含于 $\mathfrak n'$，例如幂零自同态一般不把 $W$ 映到自身之上）.}，它同样把 $r(\mathfrak g)$ 作为一个理想包含在内，并且还包含 $r(X)_s$ 与 $r(X)_n$. 因此我们只须证明 $r(\mathfrak g)=\mathfrak n'$.

让 $r(\mathfrak g)$ 经伴随作用作用于 $\mathfrak n'$，Weyl 定理表明可以找到 $r(\mathfrak g)$ 在 $\mathfrak n'$ 中的一个稳定补 $\mathfrak m$. 设 $u\in\mathfrak m$. 由于 $[r(\mathfrak g),\mathfrak n']\subset r(\mathfrak g)$，有 $[r(\mathfrak g),\mathfrak m]\subset\mathfrak m\cap r(\mathfrak g)=\{0\}$，特别地 $[r(\mathfrak g),u]=0$. 由此得出，对 $V$ 的每个不可约 $\mathfrak g$-子表示 $W$，$u|_W$ 的作用与 $\mathfrak g$ 的作用交换，因而（Schur 引理）是比例为 $\lambda_{W,u}\in K$ 的位似. 但由于 $\operatorname{tr}(u|_W)=0$，又有 $\lambda_{W,u}=0$，故 $W\subset\operatorname{Ker}(u)$. 而（再次由 Weyl 定理）$V$ 是它的不可约子表示之和. 于是 $V\subset\operatorname{Ker}(u)$，从而 $u=0$、$\mathfrak m=0$，最终 $\mathfrak n'=r(\mathfrak g)$.

\emph{第二步.}把第一步应用于伴随表示 $(\mathfrak g,\operatorname{ad}_{\mathfrak g})$. 由于态射 $\mathfrak g\xrightarrow{\ \operatorname{ad}_{\mathfrak g}\ }\mathfrak{gl}(\mathfrak g)$ 是单射，元素 $X$ 有唯一的分解 $X=X_s+X_n$，使得 $\operatorname{ad}_{\mathfrak g}(X_s)=\operatorname{ad}_{\mathfrak g}(X)_s$ 且 $\operatorname{ad}_{\mathfrak g}(X_n)=\operatorname{ad}_{\mathfrak g}(X)_n$. 显然 $[X_s,X_n]=0$.

证明对任何表示 $(V,r)$ 有 $r(X_s)=r(X)_s$ 与 $r(X_n)=r(X)_n$. 由于 $\mathfrak g$ 上的自同态 $\operatorname{ad}_{\mathfrak g}(X)$ 在过渡到商后诱导 $r(\mathfrak g)$ 上的自同态 $\operatorname{ad}_{r(\mathfrak g)}(r(X))$，有 $\operatorname{ad}_{r(\mathfrak g)}(r(X_s))=\operatorname{ad}_{r(\mathfrak g)}(r(X))_s$ 与 $\operatorname{ad}_{r(\mathfrak g)}(r(X_n))=\operatorname{ad}_{r(\mathfrak g)}(r(X))_n$. 另一方面，我们知道 $\operatorname{ad}_{\mathfrak{gl}(V)}(r(X)_s)=\operatorname{ad}_{\mathfrak{gl}(V)}(r(X))_s$ 与 $\operatorname{ad}_{\mathfrak{gl}(V)}(r(X)_n)=\operatorname{ad}_{\mathfrak{gl}(V)}(r(X))_n$，把这些自同态限制到 $r(\mathfrak g)$ 上，便得 $\operatorname{ad}_{r(\mathfrak g)}(r(X)_s)=\operatorname{ad}_{r(\mathfrak g)}(r(X))_s$ 与 $\operatorname{ad}_{r(\mathfrak g)}(r(X)_n)=\operatorname{ad}_{r(\mathfrak g)}(r(X))_n$. 由此得到
\[\operatorname{ad}_{r(\mathfrak g)}(r(X_s))=\operatorname{ad}_{r(\mathfrak g)}(r(X)_s)\quad\text{及}\quad\operatorname{ad}_{r(\mathfrak g)}(r(X_n))=\operatorname{ad}_{r(\mathfrak g)}(r(X)_n),\]
而由于 $\operatorname{ad}_{r(\mathfrak g)}$ 是单射，就得到所希望的 $r(X_s)=r(X)_s$ 与 $r(X_n)=r(X)_n$.
\end{proof}
\begin{definition*}
半单代数 $\mathfrak g$ 的元素 $X$ 称为半单的，如果 $X=X_s$；称为幂零的，如果 $X=X_n$.
\end{definition*}
具体地，定理的证明表明：$X$ 半单当且仅当 $\operatorname{ad}_X$ 可对角化，$X$ 幂零当且仅当 $\operatorname{ad}_X$ 是幂零自同态. 它还表明：对 $\mathfrak{gl}_n(K)$ 的半单子代数（例如 $\mathfrak{so}_n$、$\mathfrak{sp}_n$ 或 $\mathfrak{sl}_n$），一个元素半单（相应地，幂零）当且仅当它作为矩阵可对角化（相应地，幂零）.
\begin{exercise*}
设 $\mathfrak h$ 是半单 Lie 代数 $\mathfrak g$ 的一个 Lie 子代数，且等于它的规范化子：$\mathfrak h=\{X\in\mathfrak g,\ [X,\mathfrak h]\subset\mathfrak h\}$. 证明 $\mathfrak h$ 在 Jordan 分解下稳定.
\end{exercise*}
\end{numpara}
